\chapter{Discussion and Conclusions}\label{ch:conclusions}

\section{Answers to the research questions}\label{sec:conclusions-findings}

For the stationary Gaussian AR(1) chain under coordinatewise OU noise the
joint score is linear and known in closed form, $\score(x,t) = -\Qt x$ with
$\Qt = (e^{-2t}\Sigz + \Deltat I)^{-1}$; its couplings fill one diagonal
per order in $t$ from a tridiagonal start and return to the identity at rate
$e^{-2t}$ (RQ1). The same score is computed by sum--product on the posterior
factor graph, a tree because every observation enters through a unary
factor, and on the Gaussian chain the recursion is the Kalman filter followed
by RTS smoothing (RQ2). Off the Gaussian case the recursion stays exact but
its messages do not close, and Theorem~\ref{thm:l-lmmse} identifies Gaussian
projection as the LMMSE estimator of the covariance-matched model, so the
measured gap to a validated grid reference (median relative score error
$0.27$ at $t = 0.05$, below $10^{-2}$ by $t \approx 1$) is the price of
discarding structure beyond second moments (RQ3). In the bulk the influence
of evidence at distance $d$ decays exactly as $q(\alpha,t)^d$, and one hop of
message passing removes $85\%$ of the isolated-frame error at $t = 0.05$
(RQ4). Fisher's identity turns one forward--backward pass into the exact
likelihood gradient, so an unknown transition law is recovered from noised
sequences by EM with nothing differentiated through the recursion, the
innovation shape settling at a median of $\convshapemed$ updates against
$\convrhomed$ for the correlation (RQ5). At an equal number of clean
trajectories the recovered structure is worth $\ratiolo$--$\ratiohi\times$
lower schedule-pooled relative score error than an unstructured network and,
on a separate protocol over $\structsizes$ sizes,
$\structratiolo$--$\structratiohi\times$ against one given locality and
weight sharing, under a stated aggregation and with the arms asymmetric in
supplied information by design (RQ6).

\section{What is new, under which assumptions, and what comes next}

Three things are new: the exact structure of a chain diffusion score, with
chain BP identified with the Kalman--RTS recursion; an exact reading of what
Gaussian closure computes, which makes the measured Laplace gap the price of
a model class rather than of an implementation; and transition learning by
exact chain inference inside EM, compared at equal data, with its boundaries
mapped: no gain from added mixture capacity under the implemented grid and
cap, partial absorption of a rank-one Markov violation against failure under
long-range coupling --- near strength one tenth against the convolutional
baseline, one fifth against both --- and a two-rate convergence law.

The exact results concern one-dimensional states, linear AR(1) dynamics and
a known channel, with the locality laws proved for the Gaussian bulk only;
the non-Gaussian results are numerical, grid BP being a reference only
within its validated regime; and the learning results are statistical, with
the headline factors moving by tens of percent under defensible alternative
aggregations (Appendix~\ref{app:repro}), every capacity cell from
$C = \capcappedfrom$ upward stopping at an iteration cap, the fitted convergence
slopes being contraction factors rather than Jacobian eigenvalues, and the
misspecification scalar being an exact KL only in the Gaussian arms. Whether
generative fidelity tracks pointwise accuracy within one estimator, and
whether trained networks on real sequence data inherit any of these
properties, are untested here.

Four directions follow: rerunning the reverse-generation comparison on
convergence-selected fits rather than a fixed iteration count; a
grid-by-capacity design to separate resolution effects from overfitting; a
paired design across violation strengths and types; and, beyond the chain, the reading of
U-Nets as belief propagation (Section~\ref{sec:sm-unet}), with tree
exactness replaced by an approximation whose cost must be measured as
Chapter~\ref{ch:laplace} measures the Gaussian closure.
